
\chapter{Part I Review: From Calculus to Field Notation}

This review connects the first mathematical tools.  A derivative begins as a limit of differences; a vector begins as an object whose components depend on a basis; a Fourier coefficient begins as a projection onto a wave.  These are not separate tricks.  They are all ways of extracting one controlled component from a larger object.

\section{A guided reconstruction}

Start with a function sampled on a grid.  Differences approximate derivatives, sums approximate integrals, and dot products approximate projections.  Passing to the continuum changes the notation but not the logic:
\[
  \sum_j f_jg_j\Delta x\longrightarrow \int f(x)g(x)\dd x.
\]
Fourier analysis then says that the waves \(e^{\ii kx}\) are a basis adapted to translations.  A derivative is diagonal in that basis because \(\dd e^{\ii kx}/\dd x=\ii k e^{\ii kx}\).

\begin{exercise}
Explain in one page how projection, Fourier coefficients, and the delta function are three versions of the same idea.
\begin{answer}
Each isolates a component: projection isolates a vector component, a Fourier coefficient isolates a wave component, and the delta function isolates the value of a test function at a point.
\end{answer}
\end{exercise}
